% im Text
% \gls{<label>}		Returns defined name and marks glossary item as used.
% \Gls{<label>}		same as \gls but making the first letter upper case.
% \GLS{<label>}		same as \gls but making all letters upper case.
% plural variants:	append to any command above a "pl" to access the respective plural form.
% \Ac{<label>}, \Acs{<label>}, \Acf{<label>}, \Acl{<label>}, or plural variants: Return the respective acronym form but making the first letter upper case

% \accite{<label>}	Returns reference labeled by the cite-key, see BIM-example below.

% ngerman options:  für Dativ und Genitiv-Formen kann die long form auch im entsprechenden Kasus ausgegeben werden, sofern long-<Kasus> definiert wurden, siehe GeoLab-Beispiel unten.
% 					\acd{<label>} bzw. \acg{<label>} gibt die dativ bzw. genitiv Version des Akronyms aus. Auch hier gibt es die Version mit Großbuchstaben.


% Use {} brackets if value contains any special characters or a komma!!

\newglossaryentry{gls:mse}{name={Mean Squared Error},description={is a risk function, corresponding to the expected value of the squared error loss}}

\newglossaryentry{ci}{name={Confidence interval},description={the mean of an estimate +/- the variation in the estimate}}

\newglossaryentry{csv}{name={Comma Separated Value},description={A text file that uses a comma (,) to separate each value inputted}}